\chapter{C-Saw: Incentivizing Censorship Measurements via Circumvention}
\paperinfo{A. Nisar, A. Kashaf, I. A. Qazi, Z. A. Uzmi (LUMS, Pakistan)}{ACM SIGCOMM 2018}

\begin{keybox}[The paper in three sentences]
Censorship \emph{measurement} systems (what is blocked, where, how) struggle to find volunteers, while censorship \emph{circumvention}
tools (Tor, Lantern, VPNs) use one slow method for everything because they do not know \emph{how} each site is blocked. C-Saw
\key{combines the two}: a local proxy on the user's machine detects how each visited URL is blocked, picks the \key{cheapest
circumvention that defeats that specific blocking}, and shares what it learned through a \key{crowdsourced global database}. Users
join because pages load faster (incentive), and every user who joins makes the measurements richer.
\end{keybox}

% ------------------------------------------------------------------
\section{The problem: two kinds of systems that need each other}

\subsection{Measurement systems lack users}
Systems like OONI and Centinel want to know \emph{what is blocked, where, how and when}. They need probes inside many censored networks,
so they ask volunteers to install software or host a device. Apart from altruism there is no reason to do this, so deployment stays small.
They also test a fixed \key{target list} of URLs, which misses anything not on the list, and the lists are often outdated. And testing
URLs on someone's behalf raises an ethical question: did the user consent to their machine fetching, say, a banned political site?

\subsection{Circumvention tools are one-size-fits-all}
Tor, Lantern and VPNs send every blocked request through relays, whatever the blocking method. That always works but is slow: Tor can
increase latency by more than 5$\times$ compared with the direct path. Yet many kinds of blocking can be beaten with a cheap
\key{local fix} that needs no relay at all, if only the tool knew which kind of blocking was being used.

\subsection{How censors block, and the cheap fix for each}
A censor can interfere at each step of loading a page. Each method has a matching fix:

\begin{center}
\small
\begin{tabularx}{\linewidth}{@{}lXX@{}}
\toprule
Blocking method & What the censor does & Cheapest circumvention (local fix)\\
\midrule
\key{DNS blocking} & The ISP's resolver returns no answer, an error, or a wrong IP (e.g.\ a local host that serves a block page). &
Ask a \key{public DNS} server (e.g.\ Google DNS) instead.\\
\key{IP blocking} & Packets to the site's IP are dropped or the connection is reset. & No local fix: needs a \key{relay} (Tor, Lantern, VPN).\\
\key{HTTP blocking} & The censor reads the \texttt{Host} header and URL path of a plain HTTP GET and resets the connection, drops it, or returns a
\key{block page}. & Use \key{HTTPS}: the URL and Host header are encrypted.\\
\key{HTTPS (SNI) blocking} & The TLS handshake carries the server name in plaintext (SNI); the censor blocks on it. & \key{Domain fronting}.\\
\bottomrule
\end{tabularx}
\end{center}

\textbf{Domain fronting, explained.} A normal HTTPS connection to \texttt{youtube.com} shows the name \texttt{youtube.com} in the DNS
query (plaintext), the TLS SNI (plaintext) and the HTTP Host header (encrypted). With domain fronting, the DNS query and SNI say
\texttt{google.com} (allowed), while the encrypted Host header says \texttt{youtube.com}. The censor only sees a connection to Google. The
front-end server (Google) decrypts the request and forwards it internally to the YouTube backend. This works only when the blocked site and
the front domain are served by the same provider.

\subsection{The case study that motivated C-Saw (Pakistan)}
The authors measured from a university and homes in Pakistan, where censorship is \key{distributed}: each ISP blocks in its own way (in
centralised censorship, as in Iran, everyone sees the same blocking).
\begin{itemize}
  \item \textbf{ISP-A} blocked YouTube only at the \textbf{HTTP} level (redirect to a block page). So plain \textbf{HTTPS} gets through.
  \item \textbf{ISP-B} used \key{multi-stage blocking}: DNS blocking (resolving to a local host) \emph{and} dropping HTTP/HTTPS requests. So
        public DNS plus domain fronting gets through.
  \item Even within one ISP, different URLs were blocked by different methods.
\end{itemize}
Then they compared page load times (PLTs) for YouTube's homepage:
\begin{itemize}
  \item HTTPS/domain fronting (direct path) was much faster than going through static proxies in the US, Europe or Japan.
  \item HTTPS was much faster than \key{Tor}, whose three relays (entry, middle, exit) make the path long.
  \item For a page blocked by keyword filtering, simply putting the site's \emph{IP address} in the URL instead of its hostname was about
        1.5$\times$ faster than \key{Lantern}, which picks relays by trust, not by distance.
\end{itemize}
\textbf{Lesson}: if you know \emph{how} a URL is blocked in \emph{your} ISP, you can pick a much faster circumvention. Knowing that
requires measurement. That is the link C-Saw exploits.

% ------------------------------------------------------------------
\section{Goals and design principles}

\begin{keybox}[Goals]
\begin{itemize}
  \item \textbf{G1, scalable measurements with consent}: many users, fine-grained, and only with the user's agreement.
  \item \textbf{G2, adaptive circumvention}: pick the method with the lowest PLT for the blocking actually used.
  \item \textbf{Constraint 1}: no target list of URLs.
  \item \textbf{Constraint 2}: protect the privacy of users who report measurements.
\end{itemize}
\end{keybox}

How these are met:
\begin{itemize}
  \item \textbf{Built-in incentive}: users get faster access to blocked content, so they have a reason to install C-Saw.
  \item \textbf{Measure only what the user visits}: this removes the need for a target list \emph{and} gives consent for free, because the
        user chose to visit those URLs while browsing normally.
  \item \textbf{Detect the blocking type quickly and choose the least-overhead fix.}
\end{itemize}

\textbf{Threat model.} The censor can block, modify or reject any web connection, but will not block \emph{all} web traffic (too much
collateral damage). It may try to block access to C-Saw's database, inject false reports, or identify users who report. It cannot break
into the databases, and users trust the C-Saw proxy itself.

% ------------------------------------------------------------------
\section{Architecture}

\paperfig[0.75]{cs_fig3}{C-Saw's components: a proxy with a local database on each client; a cloud-hosted global database and server. The
client measures the direct path, circumvents when needed, sends reports of blocked URLs, and downloads the list of URLs blocked in its AS.\label{fig:csarch}}

\begin{enumerate}
  \item \key{Client-side proxy} on the user's machine. All the browser's requests go through it. It has two parts: a \key{measurement
        module} (is this URL blocked, and how?) and a \key{circumvention module} (fetch it via the direct path, a local fix, a static proxy,
        Tor, Lantern, \dots).
  \item \key{local\_DB}: an in-memory hash table on the client, one record per URL the user has visited.
  \item \key{global\_DB} plus a server (\key{server\_DB}) in the cloud: the crowdsourced list of blocked URLs from all users.
\end{enumerate}

\subsection{What is stored}
\begin{center}
\small
\begin{tabularx}{\linewidth}{@{}lX@{}}
\toprule
Field (local\_DB) & Meaning\\
\midrule
URL & the requested URL (the index)\\
AS number & the user's AS at measurement time (blocking differs per ISP)\\
$T_m$ & when it was measured\\
Status & \texttt{blocked}, \texttt{not-blocked}, or \texttt{not-measured} (never measured, or the record expired)\\
Stage-1 \dots\ Stage-$k$ blocking & the blocking seen at each stage (multi-stage, e.g.\ DNS then IP)\\
Global posted & whether this has been reported to global\_DB yet\\
\bottomrule
\end{tabularx}
\end{center}
The global\_DB has the same fields plus $T_p$ (when it was posted) and a \key{UUID} for the reporting client. It stores \emph{no IP
addresses}, so reports cannot be traced back to users. Clients report \emph{only blocked} URLs.

\subsection{Life of a client}
\begin{enumerate}
  \item \textbf{Install and register.} The user solves a ``No CAPTCHA reCAPTCHA'' (to stop bots creating thousands of fake clients). The
        server gives the client a UUID.
  \item \textbf{Download} the list of URLs that other users have reported blocked \emph{in this client's AS}.
  \item \textbf{Browse.} Every request goes through the proxy, which follows Algorithm 1 (Figure~\ref{fig:csalg}) below.
  \item \textbf{Report}, periodically, newly found blocked URLs to global\_DB, always \key{over Tor} so the censor cannot see who is reporting.
\end{enumerate}

% ------------------------------------------------------------------
\section{The measurement module}

\paperfig[0.48]{cs_alg1}{Algorithm 1: what the proxy does with each requested URL.\label{fig:csalg}}

The algorithm has three cases, depending on what the databases say about the URL:

\begin{center}
\small
\begin{tabularx}{\linewidth}{@{}lXX@{}}
\toprule
Status of URL & What the proxy does & Why\\
\midrule
\texttt{not-measured} & Send \key{redundant requests}: one on the direct path and one on a circumvention path (e.g.\ Tor) at the same
time. Measure the direct path and record the result. & We do not know if it is blocked. Detecting blocking on the direct path can be slow,
so the circumvention copy makes sure the user still gets the page fast.\\
\texttt{blocked} & Use the circumvention approach recorded for it. \key{With probability $p$}, also send a request on the direct path and
re-measure. & Get the page fast; occasionally re-check, so false reports and unblocking are caught.\\
\texttt{not-blocked} & Use only the direct path. & Fastest; no extra traffic (\key{selective redundancy}).\\
\bottomrule
\end{tabularx}
\end{center}
The result is written to local\_DB. Every record \key{expires} after $t$ seconds and becomes \texttt{not-measured} again. The proxy checks
local\_DB first, and consults its copy of the global list only if the URL is not blocked locally.

\subsection{Detecting the type of blocking}

\paperfig[0.42]{cs_fig4}{Flowchart for detecting the blocking type on the direct path. ``GDNS?'' asks whether retrying with a global (public)
DNS server succeeds.\label{fig:csdetect}}

Follow a request down the flowchart (Figure~\ref{fig:csdetect}):
\begin{enumerate}
  \item \textbf{DNS query} to the local (ISP) resolver. If there is no response or NXDOMAIN, retry with a \key{global DNS} server.
  \item \textbf{TCP connection} to the IP. If the connection gets no response or is reset: if the IP came from global DNS too, and it still
        fails, declare \key{IP blocking}; otherwise retry using global DNS.
  \item \textbf{Send HTTP/HTTPS.} No response or reset $\to$ HTTP/S blocking is suspected.
  \item \textbf{Block page detection} on the response. If it is a block page and the global-DNS address gives the same result: \key{HTTP/S
        blocking} (possibly plus DNS).
  \item If the page is correct: if it was only correct after switching to global DNS, it was \key{DNS blocking}; if it was correct straight away,
        \key{no blocking}.
\end{enumerate}
A URL is declared blocked when the circumvention path returns a response but the direct path returns nothing, or something abnormal, or a block page.

\textbf{Why detection can be slow} (average detection time per blocking type, measured in Pakistan):
\begin{center}
\begin{tabular}{@{}lr@{}}
\toprule
Blocking type & Avg.\ time to detect\\
\midrule
DNS (``server refused'') & 0.025\,s\\
HTTP (block page) & 1.8\,s\\
DNS (``server failure'') & 10.6\,s\\
TCP/IP (wait for TCP timeout) & 21\,s\\
DNS + TCP/IP (multi-stage) & 32.7\,s\\
\bottomrule
\end{tabular}
\end{center}
If the user had to wait 32.7\,s before C-Saw even tried to circumvent, nobody would use it. This is why redundant requests matter.

\subsection{Redundant requests}
For a \texttt{not-measured} URL, the proxy sends the request on the direct path \emph{and} on a circumvention path at the same time, and shows the user
whichever useful response arrives first. Two benefits:
\begin{enumerate}
  \item \textbf{Speed}: the user does not wait for detection. Sending in parallel (instead of ``try direct, detect blocking, then try Tor'') cut PLTs of
        blocked pages by \key{46--64\%}.
  \item \textbf{Telling censorship from network problems}: both paths share the user's access link. If the circumvention path works but the direct path
        does not, the problem is not the user's own connection: it is probably censorship.
\end{enumerate}
\textbf{Cost and how it is limited}: duplicates add load, which hurts pages with many embedded objects. So redundancy is \emph{selective}
(only \texttt{not-measured} URLs, and blocked URLs with probability $p$). The second copy can also be \emph{delayed} (e.g.\ by 2\,s, and not sent at
all if the direct response arrives first). HTTP POSTs are never duplicated, to avoid doing a write twice. Two copies improved median PLT by 30\%; a third
copy did not help and made the tail 17\% worse.

\subsection{Detecting block pages quickly: a two-phase algorithm}
Censors often return a normal-looking HTTP 200 page that says ``this site is blocked''. Comparing it with the real page is the reliable test, but that means
waiting for the slower circumvention copy. So:
\begin{enumerate}
  \item \textbf{Phase 1 (fast)}: look only at the direct-path response and apply an \key{HTML-tag heuristic} for block pages. If it does not look like a
        block page, show it to the user immediately.
  \item \textbf{Phase 2 (accurate)}: if it does look like one, wait for the circumvention response and \key{compare page sizes}; a block page is much smaller
        than the real page.
\end{enumerate}
On a dataset of block pages from 47 ISPs, phase 1 alone correctly classified 80\% of block pages with no false positives. If phase 1 misses a block page (false
negative), the user briefly sees it, and C-Saw refreshes the browser as soon as the real page arrives via the circumvention path.

\subsection{Choosing $p$: trust the crowd or re-check?}
\begin{center}
\begin{tabular}{@{}lll@{}}
\toprule
$p$ & Behaviour & Problem\\
\midrule
0 & always trust global\_DB's ``blocked'' & a false report forces slow circumvention forever\\
1 & always re-measure the direct path & extra traffic; global\_DB is pointless\\
$0 < p < 0.25$ & re-check occasionally & recommended balance\\
\bottomrule
\end{tabular}
\end{center}
For URLs fixed with a local fix, re-checking is free, since the local fix uses the direct path anyway.

% ------------------------------------------------------------------
\section{The circumvention module: picking a method}

\begin{keybox}[Selection rule]
\begin{enumerate}
  \item \textbf{Prefer a local fix} (public DNS for DNS blocking, HTTPS for HTTP blocking, domain fronting for HTTPS blocking): no relay, so the path is
        short.
  \item If a relay is needed (IP blocking), choose the relay method (Tor, Lantern, static proxy) with the \key{lowest moving-average PLT} for that URL.
  \item Every \key{5th} access, try a \emph{random} eligible method, in case it has become faster (exploration).
\end{enumerate}
\end{keybox}

\textbf{Example walk-through.} A user on ISP-B (DNS + HTTP/HTTPS blocking) visits a never-seen URL:
\begin{center}
\small
\begin{tabular}{@{}lll@{}}
\toprule
Request & What happens & Result\\
\midrule
1st & \texttt{not-measured}: direct + Tor in parallel; Tor's copy is shown & detect DNS + HTTPS blocking; record it\\
2nd & \texttt{blocked}: local fix = public DNS + domain fronting & fast, no relay\\
\dots, 5th & random method tried (e.g.\ Lantern) & PLT average updated\\
next & with probability $p$ also re-test direct path & detects if the site was unblocked\\
\bottomrule
\end{tabular}
\end{center}
\textbf{User preference}: a user who wants \key{anonymity} can tell C-Saw to always use an anonymous method (Tor), even when a faster local fix exists.

% ------------------------------------------------------------------
\section{Extra features}

\subsection{Tracking changes in blocking (churn)}
\begin{itemize}
  \item \textbf{Blocked $\to$ unblocked}: the record expires, the status becomes \texttt{not-measured}, redundant requests are sent, and the direct path is seen
        to work.
  \item \textbf{Unblocked $\to$ blocked}: the direct path is always used for unblocked URLs, so the blocking is noticed on the next visit.
\end{itemize}

\subsection{Keeping local\_DB small: URL aggregation}
Many devices in developing countries have little memory, so C-Saw aggregates URLs under their base URL, depending on the blocking type:
\begin{itemize}
  \item \textbf{HTTP blocking}: (a) if the base URL \texttt{foo.com/} is blocked, keep only that record; every page under it counts as blocked. (b) If only
        a derived URL \texttt{foo.com/a.html} is blocked (censors sometimes block single pages), keep a separate record for it. (c) If any URL is unblocked,
        keep one record for the base. A lookup then uses \key{longest-prefix match}.
  \item \textbf{DNS, IP, HTTPS blocking}: these block a whole hostname or IP (SNI carries only the hostname), so mark the base URL blocked, even if it was a
        derived URL that was detected.
\end{itemize}
This reduced local\_DB records by \key{55\%} for Pakistan's top-15 websites.

\subsection{Multi-homed networks}
Suppose a network sends each request at random over ISP A or ISP B, and only A blocks URL Y. The first request goes via B: ``not blocked''. The next goes via A:
``blocked''. The status keeps flipping, and C-Saw keeps re-detecting. Fix: periodically look up the ASN of the upstream provider; if more than one ASN shows up
over a short time, the network is multi-homed, and C-Saw uses the circumvention suited to the \emph{stricter} ISP for all requests to Y.

% ------------------------------------------------------------------
\section{Security and privacy}

\begin{center}
\small
\begin{tabularx}{\linewidth}{@{}lX@{}}
\toprule
Threat & Defence\\
\midrule
Censor identifies users who report & Reports go \key{over Tor}; global\_DB stores no IP addresses.\\
Censor blocks access to global\_DB & Use collectors exposed as \key{Tor hidden services} (as OONI does), so there is no well-known name to block.\\
Fake clients flood false reports & \key{reCAPTCHA} at registration limits fake identities; \key{voting} limits each client's influence (below).\\
DDoS on server/database & Host on cloud platforms with auto-scaling and DDoS protection, replicated across providers.\\
Users endangered by circumventing & Users may choose anonymous methods only; measurement is only of URLs they chose to visit (Menlo Report ethics).\\
\bottomrule
\end{tabularx}
\end{center}

\textbf{The voting scheme.} Each client $c_i$ has \emph{one} vote in total, split evenly over the $d$ URLs it reports as blocked: each gets $1/d$.
For each (URL $j$, AS $k$), the server keeps the sum of votes $s_{j,k}$ and the number of reporters $n_{j,k}$.

\textbf{Example}: an honest user reports 2 blocked URLs, giving 0.5 to each. A malicious client reports 1000 fake URLs, giving each only 0.001. A fake URL
reported only by that client has $n = 1$, $s = 0.001$: it should be trusted little. A genuinely blocked URL reported by 40 users has large $n$ \emph{and} large
$s$. A URL with large $n$ but tiny $s$ means many reporters who each report huge lists, which is suspicious too. (Inspired by PageRank.)

% ------------------------------------------------------------------
\section{Implementation and results in brief}
\begin{itemize}
  \item Built with Electron (cross-platform, works with any browser); global\_DB on MongoLab, server\_DB on Heroku; Tor and Lantern as relay options.
  \item C-Saw cut average PLT by up to \key{48\%} versus Lantern and \key{63\%} versus Tor: for a DNS-blocked page it uses a local fix instead of a relay, and for
        unblocked pages it uses the direct path, while Tor and Lantern always relay.
  \item \textbf{Pilot study}: 123 users in Pakistan over three months, 420 blocked domains, 16 ASes. Block pages were the most common method, DNS blocking second.
        It found \key{blocking of CDN servers}, which earlier studies had not seen.
  \item \textbf{Live event}: during protests on 25--28 Nov 2017, C-Saw users recorded Twitter blocked by HTTP timeout in one AS and by a block page in another,
        and Instagram blocked via DNS in several ASes: different ISPs, different methods, within minutes.
\end{itemize}

\textbf{Limitations}: coverage depends on what users visit (unpopular sites may never be measured); Tor exits can be blacklisted (but Tor is only one of the
options); users flocking to one method after another is shut down might be fingerprintable; redundant requests use more data.

% ------------------------------------------------------------------
\section{Summary comparison}
\begin{center}
\small
\begin{tabularx}{\linewidth}{@{}lXXX@{}}
\toprule
 & Measurement tools (OONI, \dots) & Circumvention tools (Tor, Lantern) & \textbf{C-Saw}\\
\midrule
Why users join & altruism & access to blocked content & \key{faster} access to blocked content\\
Which URLs & a fixed target list & n/a & \key{only what the user visits} (consent)\\
Adapts to blocking type & n/a & no, always relays & \key{yes}, local fix first\\
Shares knowledge & yes & no & yes, per-AS crowdsourced list\\
\bottomrule
\end{tabularx}
\end{center}

\begin{keybox}[Exam one-liner]
C-Saw merges censorship measurement with circumvention. A local proxy measures only URLs the user visits (consent, no target list), detects the blocking type
(DNS, IP, HTTP, HTTPS/SNI, multi-stage) with a flowchart using a global DNS retry, sends redundant requests (direct + circumvention) for unmeasured URLs to hide
slow detection (up to 32.7\,s), and uses a two-phase block-page check (HTML heuristic, then size comparison). It prefers local fixes (public DNS, HTTPS, domain
fronting) over relays, picks relays by moving-average PLT with random exploration every 5th access, and shares blocked URLs per AS via a global DB (reports over
Tor, reCAPTCHA + $1/d$ voting against fake reports). Faster pages are the incentive that makes measurement scale.
\end{keybox}

\section{Revision checklist}
\begin{itemize}
  \item Problem: measurement lacks incentive/consent and uses target lists; circumvention is one-size-fits-all and slow.
  \item Blocking types and their local fixes; domain fronting (DNS/SNI = front, encrypted Host = real site).
  \item Pakistan case study: ISPs block differently (ISP-A HTTP only; ISP-B DNS + HTTP/HTTPS); local fixes beat Tor, Lantern and static proxies.
  \item Components: proxy (measurement + circumvention modules), local\_DB, global\_DB + server\_DB; download per-AS blocked list; report over Tor.
  \item Algorithm 1: not-measured $\to$ redundant requests + measure; blocked $\to$ circumvent, re-measure with prob.\ $p$; not-blocked $\to$ direct only; records expire.
  \item Detection flowchart with GDNS retry; detection times (0.025--32.7\,s).
  \item Redundant requests: hide detection time, separate censorship from network faults; selective, optionally delayed, no POST duplication.
  \item Two-phase block-page detection; $p < 0.25$.
  \item Selection: local fix first, else min moving-average PLT, random every 5th access; anonymity preference.
  \item URL aggregation (rules per blocking type, longest-prefix match); multi-homing detection via ASN.
  \item Security: Tor reports, no IPs, hidden-service collectors, reCAPTCHA, $1/d$ voting with $s_{j,k}$ and $n_{j,k}$.
\end{itemize}
